% A content fragment, meant to be \input{} or \include{}d from a master
% document alongside "4 Human Metabolic System Model.tex" — same
% requirements (amsmath, unicode-math + fontspec, compiled with xelatex or
% lualatex for the raw Unicode symbols below).

\section{Glossary}

\subsection{Acronyms}

\begin{description}
\item[ACSM] American College of Sports Medicine.
\item[AEE] Activity energy expenditure.
\item[AMDR] Acceptable Macronutrient Distribution Range.
\item[BMI] Body mass index.
\item[BMR] Basal metabolic rate.
\item[IOM] Institute of Medicine.
\item[LBM] Lean body mass.
\item[MET] Metabolic equivalent of task.
\item[MIMO] Multi-input multi-output.
\item[SD] Sleep deprivation effect.
\item[TEF] Thermic effect of food.
\item[TEI] Total energy intake.
\item[USDA] U.S. Department of Agriculture.
\item[WHO] World Health Organization.
\end{description}

\subsection{Symbols and Variables}

\begin{description}
\item[$A$] Affine maintenance coefficient (intercept), $M(m̄) = A + B×m̄$.
\item[$AEE$] Daily activity energy expenditure (kcal/day).
\item[$B$] Affine maintenance coefficient (slope), $M(m̄) = A + B×m̄$.
\item[$BF\%$] Body-fat percentage, U.S. Navy estimate (\%).
\item[$a$] Age (years).
\item[$BMI_{des}$] Body mass index at $m_{des}$ (kg/m$^2$).
\item[$BMR$] Basal metabolic rate (kcal/day).
\item[$BMR_{adp}$] Adapted basal metabolic rate (kcal/day).
\item[$BMR_{avg}$] Average of the known basal-metabolic-rate estimates (kcal/day).
\item[$BMR_{cal}$] Calibrated basal metabolic rate (kcal/day).
\item[$BMR_{kat}$] Katch-McArdle basal metabolic rate (kcal/day).
\item[$BMR_{mif}$] Mifflin-St Jeor basal metabolic rate (kcal/day).
\item[$b$] Dietary fiber coefficient (g per 1{,}000 kcal; g/kg).
\item[$C$] Daily carbohydrate (g/day).
\item[$C_{hip}$] Hip circumference (cm).
\item[$C_{neck}$] Neck circumference (cm).
\item[$C_{waist}$] Waist circumference (cm).
\item[$c_{m}$] Boer lean-body-mass mass coefficient.
\item[$c_{h}$] Boer lean-body-mass height coefficient.
\item[$c_{0}$] Boer lean-body-mass constant term.
\item[$D$] Daily energy deficit (kcal/day).
\item[$D_{i}$] Average daily energy deficit the logs imply over calibration period $i$ (kcal/day).
\item[$D_{bmr}\%$] Daily deficit as a share of BMR (\%).
\item[$F$] Daily dietary fiber (g/day).
\item[$f$] Thermic effect of food as a fraction of $TEI$.
\item[$G$] Daily fat (g/day).
\item[$g_{liver}$] Liver glycogen store (g).
\item[$g_{musc}$] Muscle glycogen density (g/kg).
\item[$h$] Height (cm).
\item[$k$] Fat share of $TEI$ (\%).
\item[$LBM$] Lean body mass (kg).
\item[$LBM_{boer}$] Lean body mass by Boer (kg).
\item[$LBM_{navy}$] Lean body mass by U.S. Navy (kg).
\item[$L_{p}$] Calibration period length (days).
\item[$m$] Body mass (kg).
\item[$m̄$] 7-day rolling average body mass (kg).
\item[$m_{des}$] Healthy body mass (kg).
\item[$m_{fat}$] Fat mass, U.S. Navy estimate (kg).
\item[$m_{gly}$] Stored glycogen mass (g).
\item[$m_{musc}$] Skeletal muscle mass (kg).
\item[$m∞$] Equilibrium body mass (kg).
\item[$MET$] Metabolic equivalent of the activity.
\item[$n$] Number of known basal-metabolic-rate estimates in $BMR_{avg}$.
\item[$n_{p}$] Number of calibration periods.
\item[$P$] Daily protein (g/day).
\item[$p$] Protein per kg of lean body mass (g/kg/day).
\item[$q$] Carbohydrate share of $TEI$ (\%).
\item[$r$] Water bound per gram of glycogen (g/g).
\item[$s$] Sleep duration (h); skeletal-muscle fraction of $LBM$ in glycogen storage.
\item[$s_{des}$] Desire sleep duration (h).
\item[$SD$] Sleep deprivation effect (kcal/day).
\item[$t$] Elapsed time (days).
\item[$TEI$] Daily energy intake (kcal/day).
\item[$w_{i}$] Weight of calibration period $i$.
\item[$Δm$] Weekly fat-mass change (kg/week).
\item[$Δm\%$] Weekly fat-mass change relative to $m̄$ (\%).
\item[$ΔM_{i}$] Daily energy balance implied by the body-mass change over calibration period $i$ (kcal/day).
\item[$ΔM_{gly}$] Glycogen and water mass swing (kg).
\item[$γ$] Sleep deprivation rate (\%/h).
\item[$ε$] Oxygen energy yield (mL O$_2$/kcal).
\item[$η$] Sleep efficiency factor.
\item[$κ$] Oxygen uptake per MET (mL O$_2$/kg/min).
\item[$λ$] Metabolic adaptation rate (\%/week).
\item[$λt_{max}$] Metabolic adaptation ceiling (\%).
\item[$ρ$] Energy density of body fat (kcal/kg).
\item[$σ$] Sex.
\item[$τ$] Activity duration (min/day).
\end{description}

\subsection{Subscripts and Indices}

\begin{description}
\item[$adp$] Metabolic-adaptation adjusted.
\item[$avg$] Average.
\item[$boer$] Boer.
\item[$cal$] Calibrated.
\item[$des$] Desire, healthy reference value.
\item[$fat$] Fat.
\item[$gly$] Glycogen.
\item[$h$] Height term.
\item[$hip$] Hip.
\item[$i$] Calibration period index.
\item[$kat$] Katch-McArdle.
\item[$liver$] Liver.
\item[$m$] Body-mass term.
\item[$min$] Lower bound.
\item[$max$] Upper bound.
\item[$mif$] Mifflin-St Jeor.
\item[$musc$] Skeletal muscle.
\item[$navy$] U.S. Navy.
\item[$neck$] Neck.
\item[$waist$] Waist.
\item[$0$] Constant term.
\end{description}
